% P11 paper module: R23 rank-one polar drift under perfect modulus coherence.

\subsection{Rank-one polar drift under perfect modulus coherence}
\label{subsec:o3v-rankone-polar-drift}

R22 isolates the fixed-vector strong-gauge observable
\(
\mathscr G_U=(V_U-W)^*(V_U-W)
\).
The next model strengthens the R14 modulus-only firewall in a direction suggested by
the actual P11 asymptotics: even a single monotonically growing rank-one future channel
can generate a persistent polar rotation while the modulus comparison is exact.

\begin{proposition}[Canonical one-source rank-one family]
\label{prop:o3v-rankone-family}
Let \(\mathcal H\) be finite-dimensional, let \(B>0\), and let
\(e\in\mathcal H\) be a unit vector.  Put
\[
 \beta:=\langle Be,e\rangle,
 \qquad
 J:\mathbb C\to\mathcal H,
 \qquad Jz:=ze,
\]
and use the base metrics
\[
 B_R:=\beta,
 \qquad B_S:=B.
\]
The normalized base inclusion is
\[
 W=B^{1/2}J\,\beta^{-1/2},
 \qquad
 w:=We_\mathbb C=\frac{B^{1/2}e}{\sqrt\beta}.
\]
For \(t\ge1\), define
\[
 A_R(t):=t,
 \qquad
 A_S(t):=I+(t-1)ww^*,
\tag{O3V.1}\label{eq:o3v-relative}
\]
and future metrics
\[
 C_R(t):=t\beta,
 \qquad
 C_S(t):=B^{1/2}A_S(t)B^{1/2}.
\tag{O3V.2}\label{eq:o3v-future}
\]
Then the future pullback is exact,
\[
 \boxed{J^*C_S(t)J=C_R(t),}
\tag{O3V.3}\label{eq:o3v-pullback}
\]
and the modulus isometry and Jensen square-root defect satisfy
\[
 \boxed{Q_t=W,\qquad \Theta_t=0}
\tag{O3V.4}\label{eq:o3v-perfect-modulus}
\]
for every \(t\ge1\).
Moreover, if
\[
 u:=\frac{Be}{\sqrt\beta},
\]
then
\[
 \boxed{C_S(t)=B+(t-1)uu^*,}
\tag{O3V.5}\label{eq:o3v-rankone-growth}
\]
so the future target metric is a fixed positive floor plus monotone rank-one growth.
\end{proposition}

\begin{proof}
Since \(\|w\|=1\),
\[
 A_S(t)w=tw,
 \qquad
 A_S(t)^{1/2}w=\sqrt t\,w.
\]
Also
\[
 B^{1/2}w=\frac{Be}{\sqrt\beta}=u,
\]
which gives~\eqref{eq:o3v-rankone-growth}.  Furthermore
\[
 \langle u,e\rangle
 =\frac{\langle Be,e\rangle}{\sqrt\beta}
 =\sqrt\beta,
\]
so
\[
 J^*C_S(t)J
 =\beta+(t-1)|\langle u,e\rangle|^2
 =t\beta=C_R(t).
\]
The relative compression is
\[
 W^*A_S(t)W=t=A_R(t),
\]
and therefore
\[
 Q_t
 =A_S(t)^{1/2}WA_R(t)^{-1/2}
 =W.
\]
Finally
\[
 W^*A_S(t)^{1/2}W=\sqrt t=A_R(t)^{1/2},
\]
which is exactly \(\Theta_t=0\).
\end{proof}

\begin{theorem}[Rank-one polar-drift limit]
\label{thm:o3v-polar-drift}
Let \(W_t\) be the actual future normalized inclusion associated with
Proposition~\ref{prop:o3v-rankone-family}:
\[
 W_t
 :=C_S(t)^{1/2}J\,C_R(t)^{-1/2}.
\]
Then
\[
 \boxed{
 W_t1
 \longrightarrow
 \frac{Be}{\|Be\|}.
 }
\tag{O3V.6}\label{eq:o3v-future-limit}
\]
The fixed baseline direction is
\[
 W1=\frac{B^{1/2}e}{\sqrt{\langle Be,e\rangle}}.
\]
Hence
\[
 \boxed{
 W_t\longrightarrow W
 \quad\Longleftrightarrow\quad
 e\text{ is an eigenvector of }B.
 }
\tag{O3V.7}\label{eq:o3v-alignment-criterion}
\]
\end{theorem}

\begin{proof}
By~\eqref{eq:o3v-rankone-growth},
\[
 \frac1tC_S(t)
 =\frac1tB+\left(1-\frac1t\right)uu^*
 \longrightarrow uu^*
\]
in operator norm.  Norm continuity of the positive square root gives
\[
 \frac1{\sqrt t}C_S(t)^{1/2}
 =\left(\frac1tC_S(t)\right)^{1/2}
 \longrightarrow(uu^*)^{1/2}
 =\frac{uu^*}{\|u\|}.
\]
Therefore
\[
\begin{aligned}
 W_t1
 &=\frac1{\sqrt\beta}
   \left(\frac1tC_S(t)\right)^{1/2}e\\
 &\longrightarrow
 \frac1{\sqrt\beta}\frac{uu^*e}{\|u\|}
 =\frac{u}{\|u\|}
 =\frac{Be}{\|Be\|},
\end{aligned}
\]
proving~\eqref{eq:o3v-future-limit}.

If \(e\) is an eigenvector of \(B\), the two normalized directions coincide.
Conversely, equality of the two unit vectors implies
\[
 Be=cB^{1/2}e
\]
for some \(c>0\).  Multiplying by \(B^{-1/2}\) gives
\[
 B^{1/2}e=ce,
\]
so \(e\) is an eigenvector of \(B^{1/2}\), hence of \(B\).
\end{proof}

\begin{corollary}[Explicit limiting R22 angle defect]
\label{cor:o3v-angle-defect}
In the model above the source polar factor is the identity and
\(Q_t=W\), hence the actual future transport equals the polar-gauge transport
\(V_t\).  Consequently the R22 angle defect is scalar and
\[
 \boxed{
 \lim_{t\to\infty}\mathscr G_t
 =\delta_B(e)^2 I_{\mathbb C},
 }
\tag{O3V.8}\label{eq:o3v-G-limit}
\]
where
\[
 \boxed{
 \delta_B(e)^2
 =2-2\frac{\langle e,B^{3/2}e\rangle}
 {\sqrt{\langle e,B^2e\rangle\langle e,Be\rangle}}.
 }
\tag{O3V.9}\label{eq:o3v-half-full-angle}
\]
Moreover
\[
 \delta_B(e)=0
 \quad\Longleftrightarrow\quad
 e\text{ is an eigenvector of }B.
\]
\end{corollary}

\begin{proof}
The source product is
\[
 C_R(t)^{1/2}B_R^{-1/2}=\sqrt t>0,
\]
so \(U_R(t)=1\).  Equation~\eqref{eq:o3v-perfect-modulus} and the exact polar-gauge
factorization give
\[
 W_t=U_S(t)W=V_t.
\]
Thus
\[
 \mathscr G_t
 =(W_t-W)^*(W_t-W)
 =\|W_t1-W1\|^2I_{\mathbb C}.
\]
Insert Theorem~\ref{thm:o3v-polar-drift}.  Since
\[
 \|Be\|^2=\langle e,B^2e\rangle,
 \qquad
 \|B^{1/2}e\|^2=\langle e,Be\rangle,
\]
and
\[
 \langle Be,B^{1/2}e\rangle
 =\langle e,B^{3/2}e\rangle,
\]
one obtains~\eqref{eq:o3v-half-full-angle}.  The zero criterion is
Theorem~\ref{thm:o3v-polar-drift}.
\end{proof}

\begin{corollary}[Continuous strengthening of the R14 countermodel]
\label{cor:o3v-R14-strengthening}
Take
\[
 P=
 \begin{pmatrix}
 \sqrt3/2&1/2\\
 1/2&1
 \end{pmatrix},
 \qquad
 B=P^2,
 \qquad
 e=e_1.
\]
Then \(\langle Be,e\rangle=1\) and
\[
 W1=Pe_1
 =\binom{\sqrt3/2}{1/2}.
\]
For the continuous family
\[
 A_t=I+(t-1)(W1)(W1)^*,
 \qquad t\to\infty,
\]
one has \(Q_t=W\) and \(\Theta_t=0\) identically, but
\[
 \boxed{
 W_t1\longrightarrow
 \frac{
 \binom{1}{(2+\sqrt3)/4}
 }{
 \sqrt{1+((2+\sqrt3)/4)^2}
 }
 \ne W1.
 }
\tag{O3V.10}\label{eq:o3v-R14-limit}
\]
In particular
\[
 \boxed{
 \lim_{t\to\infty}\mathscr G_t
 =\left(
 2-\frac{2+5\sqrt3}{\sqrt{23+4\sqrt3}}
 \right)I_{\mathbb C}
 }
\tag{O3V.11}\label{eq:o3v-R14-angle}
\]
with
\[
 2-\frac{2+5\sqrt3}{\sqrt{23+4\sqrt3}}
 \approx0.0513796574>0.
\]
Thus no alternating sequence is needed: a monotone rank-one future blow-up with perfect
modulus coherence can converge to the wrong polar-gauge limit.
\end{corollary}

\begin{proof}
A direct multiplication gives
\[
 P^2e_1
 =\binom{1}{(2+\sqrt3)/4}.
\]
Theorem~\ref{thm:o3v-polar-drift} gives~\eqref{eq:o3v-R14-limit}; substituting the two
unit directions into Corollary~\ref{cor:o3v-angle-defect} gives
\eqref{eq:o3v-R14-angle}.
\end{proof}

\begin{remark}[R23 information firewall]
\label{rem:o3v-firewall}
Corollary~\ref{cor:o3v-R14-strengthening} is stronger than the earlier R14 alternating
countermodel in one specific sense: the future metric now has a fixed positive floor and
a single monotone rank-one channel, while the modulus comparison remains exact at every
parameter value.  Nevertheless the R22 angle defect approaches a positive constant.
Thus rank-one dominance is not by itself a polar-synchronization mechanism.

The obstruction is the mismatch between the baseline half-power direction
\[
 B^{1/2}e
\]
and the rank-one-dominant future full-power direction
\[
 Be.
\]
This model is not a counterexample to the concrete P11 metric family.  In P11 the
existing O3p/O3q results are fixed-pair or finite-block statements, and positive square
root does not commute with compression.  Therefore one must not insert a finite jet
Riesz vector into~\eqref{eq:o3v-half-full-angle} and call the result the full P11 gauge.
\end{remark}

\begin{openproblem}[R23-F: operator-level dominant-channel polar angle]
\label{open:o3v-dominant-angle}
Determine whether the concrete P11 future metric possesses a gauge-relevant
operator-level or reducing-block asymptotic of the schematic form
\[
 G_{X,U}=B_X+\tau_U r_Ur_U^*+E_U,
 \qquad \tau_U\to\infty,
\]
with remainder control strong enough to pass through the positive square root.  If so,
identify the resulting analogue of the half-power/full-power angle
\[
 \frac{B_Xr_U}{\|B_Xr_U\|}
 \quad\text{versus}\quad
 \frac{B_X^{1/2}r_U}{\|B_X^{1/2}r_U\|}
\]
and determine whether the corresponding relative \(R/S\) polar orientations
intertwine.

The present fixed-pair jet asymptotics do not supply this operator-level transfer, so no
claim about the concrete P11 limit of \(\mathscr G_U\), strong terminal transport, or a
global Object~X is made here.
\end{openproblem}
